\documentclass{article}
\usepackage{amsmath,amssymb}
\usepackage{xurl}
\usepackage[hidelinks]{hyperref}
% Shared commands for notes in docs/paper-gaps/.

\DeclareUnicodeCharacter{2080}{\ifmmode{}_{0}\else\textsubscript{0}\fi}
\DeclareUnicodeCharacter{2081}{\ifmmode{}_{1}\else\textsubscript{1}\fi}
\DeclareUnicodeCharacter{2082}{\ifmmode{}_{2}\else\textsubscript{2}\fi}
\DeclareUnicodeCharacter{2099}{\ifmmode{}_{n}\else\textsubscript{n}\fi}

\newcommand{\Lean}{\textsc{Lean}}
\ExplSyntaxOn
\NewDocumentCommand{\leanid}{m}{
  \ifmmode
    \textsf{\detokenize{#1}}
  \else
    \group_begin:
    \urlstyle{sf}
    \tl_set:Nn \l_tmpa_tl {#1}
    \tl_replace_all:Nnn \l_tmpa_tl {\_} {_}
    \exp_args:NV \path \l_tmpa_tl
    \group_end:
  \fi
}
\ExplSyntaxOff
\providecommand{\tr}{\operatorname{tr}}
\newcommand{\tnleanIssueBase}{https://github.com/LionSR/TNLean/issues/}
\newcommand{\tnleanPrBase}{https://github.com/LionSR/TNLean/pull/}
\newcommand{\tnleanDocBase}{https://sirui-lu.com/TNLean/docs/}
\newcommand{\ghissue}[1]{\href{\tnleanIssueBase#1}{GitHub issue \##1}}
\newcommand{\ghpr}[1]{\href{\tnleanPrBase#1}{PR \##1}}
\newcommand{\doclink}[1]{\href{\tnleanDocBase#1.html}{\path{#1}}}

% Dirac notation and the identity operator, as in blueprint/src/macros/common.tex.
% \providecommand keeps note-local definitions authoritative.
\usepackage{dsfont}
\providecommand{\ket}[1]{|#1\rangle}
\providecommand{\bra}[1]{\langle #1|}
\providecommand{\braket}[2]{\langle #1 | #2 \rangle}
\providecommand{\ketbra}[2]{|#1\rangle\!\langle #2|}
\providecommand{\Id}{\mathds{1}}
\providecommand{\one}{\mathds{1}}
\providecommand{\C}{\mathbb{C}}
\providecommand{\MN}[1]{M_{#1}(\C)}
\providecommand{\MD}{M_D(\C)}

% External referencing for paper sources.  The build helper writes sanitized
% auxiliary files under build/paper-aux/ and excludes duplicated source labels.
\usepackage{xr}

% Import a generated paper auxiliary file with a caller-supplied label prefix.
% The candidate paths support builds from docs/paper-gaps/, the repository root,
% and build/paper-aux/ itself.
\newcommand{\importpaperaux}[2]{%
  \IfFileExists{../../build/paper-aux/#2.aux}{%
    \externaldocument[#1]{../../build/paper-aux/#2}%
  }{%
    \IfFileExists{build/paper-aux/#2.aux}{%
      \externaldocument[#1]{build/paper-aux/#2}%
    }{%
      \IfFileExists{#2.aux}{%
        \externaldocument[#1]{#2}%
      }{}%
    }%
  }%
}

% General external reference macros
\newcommand{\paperref}[2]{\ref{#1-#2}}
\newcommand{\papereqref}[2]{(\ref{#1-#2})}

% CPSV16 source (1606.00608 / MPDO-22-12-17-2.tex) external document and shortcuts
\importpaperaux{cpsv16-}{1606.00608/MPDO-22-12-17-2-tnlean}
\newcommand{\cpsvref}[1]{\paperref{cpsv16}{#1}}
\newcommand{\cpsveqref}[1]{\papereqref{cpsv16}{#1}}

% Verdict marker: \gapnote{<kind>}{<status>}. Kind is the policy
% classification (clarification, local-correction, scope-restriction,
% unfaithful, false-source, open-gap); status is open, wip (elimination
% underway), resolved, or historical. Machine-read by the site index;
% prints a verdict line.
\newcommand{\gapnote}[2]{\par\noindent\textbf{Verdict:} #1 (#2).\par}

\title{Fibonacci String-Net: the Projector Identity at Positive Length}
\author{}
\date{2026-09-25}
\begin{document}
\maketitle
\gapnote{local-correction}{resolved}

\section{What the source prints}
Section~3.2 of \cite{Bultinck2017Anyons} writes the projector matrix product
operator on a ring of $L$ sites as $P_L=\sum_a w_aO_a^L$, where
$O_a^L=\sum_{i,j}\tr(B_a^{i_1j_1}\cdots B_a^{i_Lj_L})\ket{i_1\cdots i_L}\bra{j_1\cdots j_L}$,
and requires
\begin{align}
    P_L^2=\sum_{a,b}w_aw_bO_a^LO_b^L=\sum_aw_aO_a^L=P_L,
    \label{eq:fib_proj}
\end{align}
which ``has to hold for all $L$''.\footnote{Source traceability:
    \path{References/1511.08090/AnyonsPEPS.tex}, lines 143--156.}
No lower bound on $L$ is stated. For the Fibonacci theory, Appendix~D.1.1
gives two blocks $B_1,B_\tau$ of bond dimensions $2$ and $3$ with weights
$w_1=1/(1+\phi^2)$ and $w_\tau=\phi/(1+\phi^2)$, where
$\phi=(1+\sqrt5)/2$.\footnote{Source traceability:
    \path{References/1511.08090/AnyonsPEPS.tex}, lines 1257--1269.}

\section{The empty chain}
At $L=0$ the Hilbert space of the chain is one-dimensional and each $O_a^0$ is
the trace of the empty product, the bond dimension $\chi_a$ of the block:
$O_1^0=2$ and $O_\tau^0=3$. Using $\phi^2=\phi+1$,
\begin{align}
    P_0
     & =2w_1+3w_\tau
    =\frac{2+3\phi}{1+\phi^2}
    =\frac{2+3\phi}{2+\phi}\approx1.894,
    \label{eq:fib_empty}
\end{align}
which is neither $0$ nor $1$, so $P_0^2\neq P_0$. The fusion rules
$O_a^LO_b^L=\sum_cN_{ab}^cO_c^L$ also fail at $L=0$: they would give
$O_\tau^0O_\tau^0=9$ against $O_1^0+O_\tau^0=5$. The identity
(\ref{eq:fib_proj}) is a statement about rings with at least one site, and the
empty chain carries no physical content.

\section{The reading adopted}
The empty chain is read as a degenerate case the source does not intend, and
idempotence, like the fusion rules, is stated for every positive length
$L\geq1$.\footnote{Result:
    \leanid{FibonacciCompression.fibProjector_mul_self}.}
This is a convention, not an error in the source: at every $L\geq1$ the
statement is the source's (\ref{eq:fib_proj}) with the source's weights, and
no further hypothesis is added.

\bibliographystyle{alpha}
\bibliography{references}
\end{document}
